% diagrams.tex - all architecture diagrams of the single-cycle RV32I core in one PDF.
% Build (from this folder):   pdflatex diagrams.tex
% Every fig_*.tex also compiles on its own:   pdflatex fig_datapath.tex
\documentclass[10pt]{article}
\usepackage[a4paper, landscape, margin=1.3cm]{geometry}
\usepackage[T1]{fontenc}
\usepackage{lmodern}
\usepackage[scaled=0.92]{helvet}
\renewcommand{\familydefault}{\sfdefault}
\usepackage{tikz}
\usepackage[mode=tex]{standalone}
\usepackage{booktabs}
\usepackage{array}
\usepackage{caption}
\usepackage[hidelinks]{hyperref}
% rv32i_style.tex - shared TikZ styles for the single-cycle RV32I diagrams.
% Visual language (same as the course PDF):
%   thick border, blue fill  = holds state (clocked)
%   thin border, white fill  = combinational logic
%   green trapezium          = multiplexer (0/1 marks the input chosen by select)
%   dashed orange            = control signal;  solid navy = data
%   /32 on a wire            = bus width;       dot = wires joined
% Port names are the exact SystemVerilog port names of the module drawn.
\usetikzlibrary{arrows.meta, positioning, calc, fit, backgrounds, shapes.geometric, shapes.misc}

\definecolor{ink}{RGB}{31,56,100}
\definecolor{statefill}{RGB}{220,230,242}
\definecolor{muxdraw}{RGB}{56,118,59}
\definecolor{muxfill}{RGB}{226,239,218}
\definecolor{ctrl}{RGB}{197,90,17}
\definecolor{ctrlfill}{RGB}{252,232,216}
\definecolor{note}{RGB}{110,110,110}
\definecolor{tintIF}{RGB}{236,243,251}
\definecolor{tintID}{RGB}{240,247,235}
\definecolor{tintEX}{RGB}{255,246,229}
\definecolor{tintMEM}{RGB}{244,237,249}
\definecolor{tintWB}{RGB}{233,246,245}

\tikzset{
  >={Stealth[length=2.2mm,width=1.6mm]},
  wire/.style      = {draw=ink, line width=0.75pt},
  data/.style      = {wire, ->},
  ctl/.style       = {draw=ctrl, line width=0.75pt, dashed, dash pattern=on 3pt off 1.6pt, ->},
  ctlnoarrow/.style= {draw=ctrl, line width=0.75pt, dashed, dash pattern=on 3pt off 1.6pt},
  block/.style     = {draw=ink, line width=0.75pt, fill=white, rectangle, align=center,
                      inner sep=3pt, font=\small},
  state/.style     = {block, line width=1.8pt, fill=statefill},
  ctrlblock/.style = {block, draw=ctrl, fill=ctrlfill},
  muxst/.style     = {draw=muxdraw, fill=muxfill, line width=0.9pt},
  aluSt/.style     = {draw=ink, fill=white, line width=0.9pt},
  addr/.style      = {circle, draw=ink, fill=white, line width=0.75pt, inner sep=0.5pt,
                      minimum size=6.5mm, font=\large},
  stage/.style     = {draw=ink!70, rounded corners=6pt, line width=0.9pt},
  sig/.style       = {font=\scriptsize\ttfamily, text=ink, inner sep=1.2pt},
  csig/.style      = {font=\scriptsize\ttfamily, text=ctrl, inner sep=1.2pt},
  wlbl/.style      = {font=\tiny, text=ink, inner sep=0.8pt},
  notetxt/.style   = {font=\scriptsize\itshape, text=note, align=left},
  mlbl/.style      = {font=\tiny\bfseries, text=muxdraw, inner sep=0pt},
  title/.style     = {font=\large\bfseries, text=ink},
  dot/.style       = {circle, fill=ink, inner sep=0pt, minimum size=4pt},
  cdot/.style      = {circle, fill=ctrl, inner sep=0pt, minimum size=4pt},
  portbox/.style   = {rectangle, fill=ink, inner sep=0pt, minimum size=4pt},
  cportbox/.style  = {rectangle, fill=ctrl, inner sep=0pt, minimum size=4pt},
  flag/.style      = {sig, draw=ink, fill=white, rounded corners=1.5pt, inner sep=1.5pt},  % net label
  stub/.style      = {ctl},                                                                 % named select input
}

% Two-line block label: module name, then "instance : module".
\newcommand{\blk}[2]{\textbf{#1}\\[-1pt]{\scriptsize\ttfamily #2}}
% Same with a two-line name.
\newcommand{\blkII}[3]{\textbf{#1}\\[-2pt]\textbf{#2}\\[-1pt]{\scriptsize\ttfamily #3}}

% Bus-width mark on a horizontal wire: \wmark{(x,y)}{32}
\newcommand{\wmark}[2]{%
  \draw[draw=ink, line width=0.6pt] ($#1+(-0.07,-0.13)$) -- ($#1+(0.07,0.13)$);
  \node[wlbl, above] at ($#1+(0,0.1)$) {#2};}
% Bus-width mark on a vertical wire: \wmarkv{(x,y)}{32}
\newcommand{\wmarkv}[2]{%
  \draw[draw=ink, line width=0.6pt] ($#1+(-0.13,-0.07)$) -- ($#1+(0.13,0.07)$);
  \node[wlbl, right] at ($#1+(0.1,0)$) {#2};}

% Multiplexer, inputs on the left, output on the right.
% \muxR{name}{x}{y}{height}: defines (name-out), (name-top), (name-bot), (name-l) (left x).
\newcommand{\muxR}[4]{%
  \draw[muxst] ($(#2,#3)+(-0.28,-#4/2)$) -- ($(#2,#3)+(-0.28,#4/2)$)
               -- ($(#2,#3)+(0.28,#4/2-0.28)$) -- ($(#2,#3)+(0.28,-#4/2+0.28)$) -- cycle;
  \coordinate (#1-out) at ($(#2,#3)+(0.28,0)$);
  \coordinate (#1-top) at ($(#2,#3)+(0,#4/2-0.14)$);
  \coordinate (#1-bot) at ($(#2,#3)+(0,-#4/2+0.14)$);
  \coordinate (#1-l)   at ($(#2,#3)+(-0.28,0)$);}

% ALU shape (notched trapezium), inputs a (upper) and b (lower) on the left.
% \aluR{name}{x}{y}{height}: defines (name-a), (name-b), (name-out), (name-op) (top edge).
\newcommand{\aluR}[4]{%
  \draw[aluSt] ($(#2,#3)+(-0.65,#4/2)$) -- ($(#2,#3)+(0.65,#4/4)$) -- ($(#2,#3)+(0.65,-#4/4)$)
               -- ($(#2,#3)+(-0.65,-#4/2)$) -- ($(#2,#3)+(-0.65,-0.28)$) -- ($(#2,#3)+(-0.35,0)$)
               -- ($(#2,#3)+(-0.65,0.28)$) -- cycle;
  \coordinate (#1-a)   at ($(#2,#3)+(-0.65,#4*0.3)$);
  \coordinate (#1-b)   at ($(#2,#3)+(-0.65,-#4*0.3)$);
  \coordinate (#1-out) at ($(#2,#3)+(0.65,0)$);
  \coordinate (#1-op)  at ($(#2,#3)+(0,#4*0.375)$);
  \coordinate (#1-c)   at ($(#2,#3)+(0.12,0)$);}

% Ports on a stage boundary. \inport{x}{y}{name}{width}  (label to the left)
% The width may be empty (1-bit port); it is expanded first so \foreach variables work.
\newcommand{\portw}[1]{\edef\pwtmp{#1}\ifx\pwtmp\empty\else\,[\pwtmp]\fi}
\newcommand{\inport}[4]{%
  \node[portbox] at (#1,#2) {};
  \node[sig, anchor=east] at ($(#1,#2)+(-0.12,0)$) {#3\portw{#4}};}
\newcommand{\cinport}[4]{%
  \node[cportbox] at (#1,#2) {};
  \node[csig, anchor=east] at ($(#1,#2)+(-0.12,0)$) {#3\portw{#4}};}
\newcommand{\outport}[4]{%
  \node[portbox] at (#1,#2) {};
  \node[sig, anchor=west] at ($(#1,#2)+(0.12,0)$) {#3\portw{#4}};}
\newcommand{\coutport}[4]{%
  \node[cportbox] at (#1,#2) {};
  \node[csig, anchor=west] at ($(#1,#2)+(0.12,0)$) {#3\portw{#4}};}
% Control input port on the top edge (label above).
\newcommand{\ctopport}[4]{%
  \node[cportbox] at (#1,#2) {};
  \node[csig, anchor=south] at ($(#1,#2)+(0,0.1)$) {#3\portw{#4}};}
% Data input port on the top edge (label above).
\newcommand{\topport}[4]{%
  \node[portbox] at (#1,#2) {};
  \node[sig, anchor=south] at ($(#1,#2)+(0,0.1)$) {#3\portw{#4}};}

% Legend strip. \legend{x}{y}
\newcommand{\legend}[2]{%
  \begin{scope}[shift={(#1,#2)}, font=\scriptsize]
    \node[state, minimum width=6mm, minimum height=3.5mm, inner sep=0pt] at (0,0) {};
    \node[anchor=west] at (0.45,0) {holds state};
    \node[block, minimum width=6mm, minimum height=3.5mm, inner sep=0pt] at (3.2,0) {};
    \node[anchor=west] at (3.65,0) {combinational};
    \draw[muxst] (6.2,-0.25) -- (6.2,0.25) -- (6.5,0.1) -- (6.5,-0.1) -- cycle;
    \node[anchor=west] at (6.6,0) {multiplexer};
    \draw[data] (8.7,0) -- (9.4,0); \node[anchor=west] at (9.45,0) {data};
    \draw[ctl] (10.4,0) -- (11.1,0); \node[anchor=west] at (11.15,0) {control};
    \node[portbox] at (12.6,0) {}; \node[anchor=west] at (12.7,0) {module port};
    \draw[wire] (14.6,0) -- (15.3,0); \draw[draw=ink, line width=0.6pt] (14.88,-0.13) -- (15.02,0.13);
    \node[wlbl] at (14.95,0.24) {32}; \node[anchor=west] at (15.35,0) {bus width};
  \end{scope}}

\captionsetup{font=small, labelfont=bf}
\newcommand{\sv}[1]{\texttt{#1}}
\newcolumntype{L}[1]{>{\raggedright\arraybackslash}p{#1}}

\begin{document}

\begin{center}
  {\LARGE\bfseries\color{ink} Single-cycle RV32I core --- architecture diagrams}\\[4pt]
  {\large Iteration 1 of the pipelined RV32I + INT8 SIMD project \quad$\cdot$\quad RTL in \sv{single\_cycle/rtl/}}
\end{center}

\noindent The core executes one instruction per clock cycle. It is built from five
\emph{stage modules} (\sv{if\_stage}, \sv{id\_stage}, \sv{ex\_stage}, \sv{mem\_stage},
\sv{wb\_stage}); \sv{riscv\_core} only wires them together. In iteration 2 the pipeline
registers IF/ID, ID/EX, EX/MEM and MEM/WB go on exactly the wires that cross the stage
borders in Figure~\ref{fig:dp}, and the stage modules stay as they are.
Every name in the figures is the exact SystemVerilog port or signal name.

\medskip
\noindent\textbf{How to read the figures.} Thick border: holds state. Thin border: combinational.
Green trapezium: multiplexer (the numbers are the select values). Solid navy: data.
Dashed orange: control. A small square on a stage border is a module port; \texttt{/32} is a bus width;
a dot joins wires. In Figure~\ref{fig:dp}, rounded flags such as \fbox{\sv{rs2\_data}} are net labels
(the same wire as the signal of that name), and each select is drawn as a named dashed stub.

\medskip
\noindent\begin{minipage}{\linewidth}\centering
  \includestandalone[width=0.9\linewidth]{fig_datapath}
  \captionof{figure}{The whole single-cycle datapath (\sv{rtl/riscv\_core.sv}), with the five stage modules shaded.
           \sv{redirect\_pc} / \sv{pc\_redirect} run from EX back to IF (next PC);
           \sv{wb\_data} runs from WB back to ID (register-file write port).}
  \label{fig:dp}
\end{minipage}

% ------------------------------------------------------------------
\clearpage
\section*{Control unit (\sv{rtl/control.sv}, PDF Table 9.2)}
A dash means ``does not matter''; the code then drives the safe default (0 / \sv{ALU\_ADD} / \sv{IMM\_I} / \sv{WB\_ALU}).
\sv{reg\_write} and \sv{mem\_write} are never don't-cares. \sv{funct3} and \sv{instr[30]} (\sv{funct7\_5}) refine the ALU operation
for \sv{OP\_REG} / \sv{OP\_IMM}; \sv{instr[30]} selects \sv{sub} \emph{only} in \sv{OP\_REG}, and \sv{sra}/\sv{srai} in both.

\medskip
\begin{center}
\small
\begin{tabular}{l l c c c l c c c c c c c}
\toprule
group & opcode & \sv{reg\_} & \sv{alu\_src} & \sv{alu\_src} & \sv{alu\_op} & \sv{imm\_} & \sv{mem\_} & \sv{mem\_} & \sv{wb\_} & \sv{branch} & \sv{jump} & \sv{jalr}\\
      &        & \sv{write} & \sv{\_a\_pc} & \sv{\_b\_imm} &        & \sv{sel}  & \sv{read} & \sv{write} & \sv{sel} &  &  & \\
\midrule
add sub \dots\ and      & \sv{0110011} OP\_REG    & 1 & 0 & 0 & f3 / f7  & --  & 0 & 0 & ALU & 0 & 0 & 0\\
addi slti \dots\ srai   & \sv{0010011} OP\_IMM    & 1 & 0 & 1 & f3       & I   & 0 & 0 & ALU & 0 & 0 & 0\\
lb lh lw lbu lhu        & \sv{0000011} OP\_LOAD   & 1 & 0 & 1 & ADD      & I   & 1 & 0 & MEM & 0 & 0 & 0\\
sb sh sw                & \sv{0100011} OP\_STORE  & 0 & 0 & 1 & ADD      & S   & 0 & 1 & --  & 0 & 0 & 0\\
beq \dots\ bgeu         & \sv{1100011} OP\_BRANCH & 0 & --& --& --       & B   & 0 & 0 & --  & 1 & 0 & 0\\
lui                     & \sv{0110111} OP\_LUI    & 1 & --& 1 & PASSB    & U   & 0 & 0 & ALU & 0 & 0 & 0\\
auipc                   & \sv{0010111} OP\_AUIPC  & 1 & 1 & 1 & ADD      & U   & 0 & 0 & ALU & 0 & 0 & 0\\
jal                     & \sv{1101111} OP\_JAL    & 1 & --& --& --       & J   & 0 & 0 & PC4 & 0 & 1 & 0\\
jalr                    & \sv{1100111} OP\_JALR   & 1 & 0 & 1 & ADD      & I   & 0 & 0 & PC4 & 0 & 0 & 1\\
fence                   & \sv{0001111} OP\_FENCE  & 0 & --& --& --       & --  & 0 & 0 & --  & 0 & 0 & 0\\
ecall ebreak            & \sv{1110011} OP\_SYSTEM & 0 & --& --& --       & --  & 0 & 0 & --  & 0 & 0 & 0\\
anything else           & ---                     & 0 & --& --& --       & --  & 0 & 0 & --  & 0 & 0 & 0\\
\bottomrule
\end{tabular}
\end{center}
\noindent Extra outputs: \sv{system} = 1 for \sv{OP\_SYSTEM} (the testbench stops on \sv{ebreak}),
\sv{illegal} = 1 for any other unknown opcode (including custom-0, \sv{0001011}, reserved for the SIMD extension).

\bigskip
\noindent\begin{tabular}{@{}L{0.48\linewidth} L{0.48\linewidth}@{}}
\textbf{ALU operations} (\sv{riscv\_pkg::ALU\_*})\par\smallskip
\small
\begin{tabular}{r l l}
\toprule
code & name & result\\ \midrule
0 & \sv{ALU\_ADD}  & a + b\\
1 & \sv{ALU\_SUB}  & a -- b\\
2 & \sv{ALU\_SLL}  & a $\ll$ b[4:0]\\
3 & \sv{ALU\_SLT}  & signed a $<$ b\\
4 & \sv{ALU\_SLTU} & unsigned a $<$ b\\
5 & \sv{ALU\_XOR}  & a \^{} b\\
6 & \sv{ALU\_SRL}  & a $\gg$ b[4:0] (zeros in)\\
7 & \sv{ALU\_SRA}  & a $\gg$ b[4:0] (sign in)\\
8 & \sv{ALU\_OR}   & a | b\\
9 & \sv{ALU\_AND}  & a \& b\\
10& \sv{ALU\_PASSB}& b (for \sv{lui})\\ \bottomrule
\end{tabular}
&
\textbf{Immediate formats} (\sv{riscv\_pkg::IMM\_*})\par\smallskip
\small
\begin{tabular}{l l}
\toprule
sel & value\\ \midrule
\sv{IMM\_I} & sext(\sv{instr[31:20]})\\
\sv{IMM\_S} & sext(\{\sv{instr[31:25]}, \sv{instr[11:7]}\})\\
\sv{IMM\_B} & sext(\{\sv{[31]}, \sv{[7]}, \sv{[30:25]}, \sv{[11:8]}, 0\})\\
\sv{IMM\_U} & \{\sv{instr[31:12]}, 12'b0\}\\
\sv{IMM\_J} & sext(\{\sv{[31]}, \sv{[19:12]}, \sv{[20]}, \sv{[30:21]}, 0\})\\ \bottomrule
\end{tabular}

\bigskip
\textbf{Write-back select} (\sv{riscv\_pkg::WB\_*}): \sv{WB\_ALU} = 0, \sv{WB\_MEM} = 1, \sv{WB\_PC4} = 2.
\end{tabular}

% ------------------------------------------------------------------
\clearpage
\section*{The five stage modules}
Port lists are the exact module ports. The ``connects to'' column is the wiring in \sv{riscv\_core.sv};
in the pipeline, every signal that crosses into the next stage is stored in that stage's pipeline register.

\begin{figure}[h]
  \centering
  \includestandalone[width=0.62\linewidth]{fig_if_stage}
  \caption{\sv{if\_stage} (\sv{rtl/if\_stage.sv}): PC register, instruction memory, \sv{pc+4}, next-PC mux.}
\end{figure}
\begin{center}\small
\begin{tabular}{l c c L{11cm}}
\toprule
port & dir & width & connects to\\ \midrule
\sv{clk}, \sv{rst} & in & 1 & core clock, synchronous reset (\sv{pc} $\leftarrow$ \sv{RESET\_ADDR})\\
\sv{pc\_redirect} & in & 1 & \sv{ex\_stage.pc\_redirect}\\
\sv{redirect\_pc} & in & 32 & \sv{ex\_stage.redirect\_pc}\\
\sv{pc} & out & 32 & \sv{ex\_stage.pc}, \sv{dbg\_pc}\\
\sv{pc\_plus4} & out & 32 & \sv{wb\_stage.pc\_plus4}\\
\sv{instr} & out & 32 & \sv{id\_stage.instr}, \sv{dbg\_instr}\\ \bottomrule
\end{tabular}
\end{center}

\clearpage
\noindent\begin{minipage}[c]{0.5\linewidth}\centering
  \includestandalone[width=\linewidth]{fig_id_stage}
  \captionof{figure}{\sv{id\_stage} (\sv{rtl/id\_stage.sv}): field slicing, control unit, register file, immediate generator.}
\end{minipage}\hfill
\begin{minipage}[c]{0.47\linewidth}\small
\begin{tabular}{L{4.2cm} c c L{5.5cm}}
\toprule
port & dir & width & connects to\\ \midrule
\sv{clk} & in & 1 & core clock (register-file write)\\
\sv{instr} & in & 32 & \sv{if\_stage.instr}\\
\sv{rf\_we} & in & 1 & \sv{reg\_write \& \textasciitilde rst}\\
\sv{rf\_waddr} & in & 5 & \sv{rd\_addr} (this stage's own output)\\
\sv{rf\_wdata} & in & 32 & \sv{wb\_stage.wb\_data}\\
\sv{rs1\_addr}, \sv{rs2\_addr} & out & 5 & not used yet; the pipeline's forwarding and hazard units need them\\
\sv{rd\_addr} & out & 5 & \sv{rf\_waddr}\\
\sv{funct3} & out & 3 & \sv{ex\_stage.funct3} (branch condition), \sv{mem\_stage.funct3} (access size)\\
\sv{rs1\_data} & out & 32 & \sv{ex\_stage.rs1\_data}\\
\sv{rs2\_data} & out & 32 & \sv{ex\_stage.rs2\_data}, \sv{mem\_stage.store\_data}\\
\sv{imm} & out & 32 & \sv{ex\_stage.imm}\\
\sv{reg\_write} & out & 1 & \sv{rf\_we} (gated with \sv{\textasciitilde rst})\\
\sv{alu\_src\_a\_pc}, \sv{alu\_src\_b\_imm}, \sv{alu\_op}[3:0] & out & & \sv{ex\_stage} operand muxes and ALU\\
\sv{branch}, \sv{jump}, \sv{jalr} & out & 1 & \sv{ex\_stage} redirect decision\\
\sv{mem\_read} & out & 1 & not used yet; the pipeline's load-use stall needs it\\
\sv{mem\_write} & out & 1 & \sv{mem\_stage.mem\_write} (gated with \sv{\textasciitilde rst})\\
\sv{wb\_sel} & out & 2 & \sv{wb\_stage.wb\_sel}\\
\sv{illegal}, \sv{system} & out & 1 & \sv{dbg\_illegal}, \sv{dbg\_system} (testbench stop conditions)\\ \bottomrule
\end{tabular}
\end{minipage}

\clearpage
\begin{figure}[h]
  \centering
  \includestandalone[width=0.66\linewidth]{fig_ex_stage}
  \caption{\sv{ex\_stage} (\sv{rtl/ex\_stage.sv}): operand muxes, ALU, branch unit, target adders and the redirect decision.}
\end{figure}
\begin{center}\small
\begin{tabular}{l c c L{11cm}}
\toprule
port & dir & width & connects to\\ \midrule
\sv{pc}, \sv{rs1\_data}, \sv{rs2\_data}, \sv{imm} & in & 32 & from IF / ID\\
\sv{alu\_src\_a\_pc}, \sv{alu\_src\_b\_imm}, \sv{alu\_op}, \sv{funct3}, \sv{branch}, \sv{jump}, \sv{jalr} & in & & from ID (control, \sv{funct3})\\
\sv{alu\_y} & out & 32 & \sv{mem\_stage.addr}, \sv{wb\_stage.alu\_y}\\
\sv{take\_branch} & out & 1 & observation only\\
\sv{pc\_redirect}, \sv{redirect\_pc} & out & 1, 32 & \sv{if\_stage} (next-PC mux); the pipeline's flush signal later\\ \bottomrule
\end{tabular}
\end{center}

\clearpage
\begin{figure}[h]
  \centering
  \includestandalone[width=0.62\linewidth]{fig_mem_stage}
  \caption{\sv{mem\_stage} (\sv{rtl/mem\_stage.sv}): the load/store unit around the data memory.}
\end{figure}
\begin{center}\small
\begin{tabular}{l c c L{11cm}}
\toprule
port & dir & width & connects to\\ \midrule
\sv{clk} & in & 1 & core clock (data-memory write)\\
\sv{addr} & in & 32 & \sv{ex\_stage.alu\_y}\\
\sv{store\_data} & in & 32 & \sv{id\_stage.rs2\_data}\\
\sv{funct3} & in & 3 & \sv{id\_stage.funct3}\\
\sv{mem\_write} & in & 1 & \sv{mem\_write \& \textasciitilde rst}\\
\sv{load\_data} & out & 32 & \sv{wb\_stage.load\_data}\\ \bottomrule
\end{tabular}
\end{center}

\bigskip
\begin{figure}[h]
  \centering
  \includestandalone[width=0.42\linewidth]{fig_wb_stage}
  \caption{\sv{wb\_stage} (\sv{rtl/wb\_stage.sv}): the write-back mux; \sv{wb\_data} returns to \sv{id\_stage.rf\_wdata}.}
\end{figure}

\end{document}
